\chapter{Cómo se verificó la placa}\label{cap:verif}

Una placa puede verse perfecta en pantalla y tener una pista a \SI{0.05}{\milli\meter} de otra, un pin sin conectar
o una resistencia girada al revés en el archivo de montaje. Ninguno de esos errores se ve a simple vista y todos
cuestan dinero y semanas cuando la placa llega. Por eso, en este proyecto cada paso del ruteo terminó con una ronda de
verificaciones hechas por herramientas, no ``a ojo''.

En este capítulo vas a ver:
\begin{itemize}
  \item qué revisa cada herramienta (DRC, paridad, ERC, \archivo{verify.py}) y por qué una sola no alcanza;
  \item el chequeo propio \archivo{revisar_pistas.py}, con sus chequeos A a M uno por uno;
  \item las herramientas de medición y de cruces entre capas;
  \item cómo se revisaron los Gerber sin escribir un lector propio;
  \item cómo se corrigió y comprobó el archivo de posiciones (CPL) para el montaje;
  \item qué hay en la carpeta \archivo{fabricacion/} y qué archivo se sube a JLCPCB;
  \item la tabla con el estado final verificado.
\end{itemize}

\begin{idea}[tres preguntas distintas]
Para que una placa esté bien hay que responder tres preguntas diferentes, y cada herramienta responde una:
\begin{enumerate}
  \item \textbf{¿El circuito está bien dibujado?} Pines de salida unidos entre sí, entradas al aire, alimentaciones
  sin fuente\ldots{} Eso lo revisa el \textbf{ERC} sobre el esquemático.
  \item \textbf{¿La placa es el mismo circuito que el esquemático?} Una pista que une dos redes que no debían unirse
  o una pieza que no está en el esquemático. Eso lo revisa la \textbf{paridad}.
  \item \textbf{¿La placa se puede fabricar y cumple las reglas?} Separaciones, anchos, anillos, bordes. Eso lo revisa
  el \textbf{DRC}.
\end{enumerate}
Y queda una cuarta que ninguna herramienta de KiCad responde: \textbf{¿la placa sigue nuestras reglas de buen
diseño?} (entrar al pad por el centro, desacoplo pegado al pin, nada crítico sobre una ranura\ldots). Para eso se
escribió \archivo{revisar_pistas.py}.
\end{idea}

\section{DRC con paridad}\label{sec:c11-drc}

El DRC (\emph{Design Rule Check}) de KiCad compara la geometría de la placa con las reglas configuradas: las
restricciones de la placa (los mínimos de JLCPCB), las clases de red (anchos y separaciones de diseño) y las reglas
especiales de \archivo{pcb.kicad_dru}. Esas reglas se explican en el capítulo~\ref{cap:jlc}.

Se corre por línea de comandos, siempre igual:
\begin{center}
\small\ttfamily kicad-cli pcb drc --schematic-parity --severity-all --refill-zones --save-board pcb.kicad\_pcb
\end{center}
Cada opción tiene una razón:
\begin{itemize}
  \item \texttt{--schematic-parity}: además de la geometría, compara la placa con el esquemático (paridad);
  \item \texttt{--severity-all}: informa errores \emph{y} advertencias; una advertencia escondida es un error que
  nadie mira;
  \item \texttt{--refill-zones}: vuelve a rellenar las zonas de cobre antes de revisar. Sin esto el DRC revisaría el
  relleno viejo, que puede no coincidir con las pistas nuevas;
  \item \texttt{--save-board}: guarda ese relleno en la placa. Además es lo que genera los 440 \emph{teardrops}
  (capítulo~\ref{cap:ruteo}) sin abrir la interfaz gráfica.
\end{itemize}

\begin{decision}[ningún chequeo en ``ignorar'']
\textbf{Qué se decidió:} ninguna regla del DRC quedó en ``ignorar'' y ninguna severidad se bajó para que pasara.\par
\textbf{Alternativas:} dejar las cinco reglas que KiCad trae ignoradas de fábrica, o excluir violaciones puntuales
(KiCad permite marcarlas como ``excluidas'').\par
\textbf{Por qué:} una regla ignorada no avisa nunca. \archivo{tools/config_reglas.py} pasó a advertencia las cinco
que KiCad trae ignoradas: \texttt{footprint\_filters\_mismatch}, \texttt{footprint\_type\_mismatch},
\texttt{missing\_courtyard}, \texttt{track\_not\_centered\_on\_via} y
\texttt{tuning\_profile\_track\_geometries}.\par
\textbf{Fuente:} pedido del dueño; \archivo{PLAN_RUTEO.md} §17.
\end{decision}

Al activarlas aparecieron 7 avisos de \texttt{footprint\_filters\_mismatch}. El aviso dice que el símbolo del
esquemático tiene una lista de footprints aceptados y que el footprint que se usó no está en esa lista. Pasaba con
símbolos genéricos de KiCad (\texttt{Conn\_01x02}, \texttt{Jumper\_2\_Open} y el símbolo del TLV9062) cuyo filtro no
incluía el footprint real (las borneras, la tira de pines y el TSOT-23-8 de pads de \SI{0.40}{\milli\meter}).

\begin{ojo}[arreglar de fondo, no tapar]
La salida fácil era excluir los 7 avisos. Se hizo lo contrario: \archivo{tools/buildlib.py} agrega a la
biblioteca del proyecto copias de esos tres símbolos con el filtro correcto, y \archivo{tools/design.py} apunta J1,
J5, J7, JP1, JP2, U3 y U9 a esas copias. Un primer intento cambió el filtro de otra forma y produjo un error nuevo en
el ERC (\texttt{lib\_symbol\_mismatch}); por eso se rehízo así. Al regenerar, el esquemático solo cambió en esos
\texttt{lib\_id}.
\end{ojo}

\textbf{Resultado final:} DRC con 0 violaciones (errores y advertencias) y paridad con 0 diferencias, con todas las
reglas activas.

\section{ERC del esquemático}\label{sec:c11-erc}

El ERC (\emph{Electrical Rule Check}) revisa el esquemático: cada pin tiene un tipo (entrada, salida, pasivo,
alimentación\ldots) y el ERC avisa de combinaciones raras, como dos salidas unidas o una red de alimentación que nadie
alimenta. Se corre con \texttt{kicad-cli sch erc --severity-all}. Resultado: \textbf{0} con las severidades del
proyecto.

KiCad trae 4 chequeos de ERC ignorados de fábrica. Se activaron en una copia para ver qué decían: solo aparecen 3
avisos \texttt{four\_way\_junction}, que son uniones de un pin con tres cables (en el cátodo de D1 y en los símbolos de
alimentación \texttt{\#PWR006} y \texttt{\#PWR018}). Eléctricamente son correctas; el aviso existe porque en un
esquemático impreso un cruce de cuatro cables se confunde con un cruce sin unión. El dueño decidió dejarlos como
están (capítulo~\ref{cap:banco}).

\section{\texorpdfstring{\archivo{verify.py}}{verify.py}: el esquemático contra el diseño}\label{sec:c11-verify}

El esquemático de KiCad no se dibujó a mano: lo genera \archivo{tools/gen_sch_wired.py} a partir de
\archivo{tools/design.py}, que es la ``fuente de verdad'' (cada pieza, su valor, su footprint y a qué red va cada
pin). \archivo{tools/verify.py} cierra el círculo: exporta la lista de redes del esquemático
(\texttt{kicad-cli sch export netlist}) y la compara con \archivo{tools/target.json}, que sale de
\archivo{design.py}. Para cada pieza revisa que:
\begin{itemize}
  \item exista con el mismo valor y el mismo footprint;
  \item cada pin esté en una red cuyos miembros son exactamente los esperados, y que el nombre de la red coincida
  (las redes sin etiqueta pueden tener el nombre automático \texttt{Net-(\ldots)});
  \item los pines que deben quedar libres estén solos.
\end{itemize}
Resultado: \textbf{127 piezas, 68 redes, 0 problemas}. Junto con la paridad, esto garantiza la cadena completa:
diseño $\to$ esquemático $\to$ placa.

\section{El chequeo propio: \texorpdfstring{\archivo{revisar_pistas.py}}{revisar\_pistas.py}}\label{sec:c11-revisar}

El DRC de KiCad sabe de separaciones y anchos, pero no sabe que un condensador de desacoplo debe estar pegado a su
pin, o que la entrada de un op-amp no debe pasar sobre un hueco del plano. Esas reglas salen de los datasheets y de las
notas de aplicación (capítulos~\ref{cap:capas} y~\ref{cap:ruteo}). \archivo{tools/ruteo/revisar_pistas.py} las
revisa leyendo la placa con la API de KiCad. Sale con código 1 si encuentra algo, así que no se puede ``pasar por
alto''.

Cada chequeo tiene una letra. Esto es lo que revisa cada uno y por qué:

\begin{description}[style=nextline,leftmargin=1.2em]
  \item[A. Orientación a 45°]
  Cada segmento de pista debe ir a un múltiplo de \SI{45}{\degree} (horizontal, vertical o diagonal). Es la regla de
  estilo del proyecto: las pistas en ángulos libres son difíciles de revisar y de mantener.

  \item[B. Ángulo entre segmentos de al menos 135°]
  Dos segmentos unidos no pueden formar un ángulo de \SI{90}{\degree} ni agudo. A frecuencias de audio una esquina
  a \SI{90}{\degree} no daña la señal; es la regla de estilo que pidió el dueño para todo el proyecto. El DRC
  también lo revisa (\texttt{track\_angle}). El script agrega un caso fino: cuando dos pistas se unen en el centro de
  un pad o de una vía, calcula la esquina interior entre sus bordes; si cae dentro del cobre del pad, no hay esquina
  de pista y se lista como ``cumple'' (70 casos en la placa final).

  \item[C. Entrada al pad por el centro y por su eje]
  La pista debe terminar en el centro del pad y llegar paralela a uno de sus ejes. Si llega torcida, el pad queda
  cargado de un lado y la pieza puede ``caminar'' al fundirse la soldadura. Hay 8 puntos aceptados por el dueño: los
  dos extremos del enlace del nodo \net{SUM} en cada canal (capítulo~\ref{cap:banco}).

  \item[D. Segmentos de menos de 0,1 mm]
  Un segmento tan corto casi siempre es un residuo de una edición y crea un quiebre invisible.

  \item[E. Extremos sueltos y pistas duplicadas]
  Un \emph{stub} (pista que no lleva a ningún lado) es una antena pequeña y cobre que sobra. Las pistas duplicadas o
  superpuestas confunden la medición y el DRC.

  \item[F. Vías sin usar]
  Una vía que conecta menos de dos capas (por pista o por zona) no sirve para nada y ocupa lugar.

  \item[G. Política de capas]
  Con 4 capas, \capa{In1.Cu} no lleva pistas: es el plano de tierra entero. Las redes críticas (clases Analogica,
  Reloj y Conmutado) solo van en \capa{F.Cu}, salvo \net{DC_MON}/\net{DC_MON2} (un nivel de DC) y las \net{IEPE_FEED*}
  (riel de DC). \capa{In2.Cu} y \capa{B.Cu} quedan para alimentación y señales lentas.

  \item[H. Cruces sobre puentes (de la época de 2 capas)]
  Con 2 capas, la referencia de \capa{F.Cu} era \capa{B.Cu} y este chequeo vigilaba que nada crítico pasara sobre una
  pista de \capa{B.Cu}. Con 4 capas la referencia es \capa{In1.Cu}, así que este chequeo ya no aplica.

  \item[I. Plano continuo bajo cada pista crítica]
  Para cada pista crítica de \capa{F.Cu} revisa, punto por punto, que debajo haya cobre del relleno de tierra de
  \capa{In1.Cu}. Es la regla de la corriente de retorno (capítulo~\ref{cap:capas}). Alrededor de cada agujero que no
  es de tierra el plano tiene un anillo libre; si la pista termina justo en ese pad o vía, esos puntos se listan como
  información (436 en la placa final), porque ese anillo existe siempre.

  \item[J. Analógica paralela a digital]
  Una señal analógica no puede correr paralela a una de reloj o digital a menos de \SI{1}{\milli\meter} por más de
  \SI{2}{\milli\meter}, en la misma capa. Dos pistas paralelas forman un condensador y un acople inductivo: cuanto más
  largo el tramo, más ruido pasa. Al principio solo miraba \capa{F.Cu}; después se extendió a todas las capas.

  \item[K. Vía de tierra junto al pad]
  Cada condensador de desacoplo y cada pin de tierra SMD de un integrado debe tener una vía de tierra a
  $\le\SI{1.0}{\milli\meter}$ (borde a borde) de su pad. Sin esa vía, la corriente de alta frecuencia tiene que viajar
  por una pista hasta el plano y el desacoplo pierde efecto. En la placa final cumplen 46 pads.

  \item[L. Desacoplo pegado a su pin]
  Cada condensador de desacoplo debe quedar a $\le\SI{1.0}{\milli\meter}$ (borde a borde) del pin que desacopla. Hay
  una tabla (\texttt{DESACOPLO}) que dice qué condensador va con qué pin. Para los del MCP4725 se usa el límite de su
  datasheet, \SI{4}{\milli\meter} \doc{DS22039D §8.3}. Los que quedaron más lejos y el dueño aceptó se listan aparte
  con su motivo.

  \item[M. Serigrafía sobre una vía]
  Las vías van tapadas (\emph{plugged}), así que no tienen apertura en la máscara y el DRC no ve un problema si la
  serigrafía pasa por encima. Pero la tinta sobre el relieve de una vía tapada se imprime mal. Este chequeo mide la
  distancia de cada vía al trazo real de la serigrafía (con \archivo{geom_serigrafia.py}), no a su caja envolvente. La
  primera vez encontró 60 casos; se corrigieron todos.
\end{description}

\begin{table}[htbp]
\centering
\caption{Resultado de \archivo{revisar_pistas.py} sobre la placa final \medido.}\label{tab:c11-revisar}
\small
\begin{tabular}{@{}lr@{}}
\toprule
Grupo & Cantidad\\\midrule
\textbf{Hallazgos} (cualquier chequeo) & \textbf{0}\\
B: uniones con la esquina dentro del pad (cumple) & 70\\
C: entradas aceptadas por el dueño (nodo SUM) & 8\\
I: puntos en el anillo libre de su propio pad o vía (información) & 436\\
K: pads con su vía de tierra a $\le\SI{1}{\milli\meter}$ & 46\\
L: desacoplos a $\le\SI{1}{\milli\meter}$ & 9\\
L: desacoplos dentro del límite del datasheet (MCP4725) & 4\\
L: desacoplos aceptados por el dueño & 11\\
\bottomrule
\end{tabular}
\end{table}

\begin{ojo}[cambiar un chequeo también es una decisión]
Cada vez que se cambió la lógica de un chequeo (por ejemplo, que el anillo libre de un pad propio cuente como
información en el I, o el límite de \SI{4}{\milli\meter} para el MCP4725 en el L), el cambio se informó al dueño en
el resumen de ese bloque y queda documentado en \archivo{PLAN_RUTEO.md}. Un chequeo que se afloja en silencio es
igual que una regla ignorada.
\end{ojo}

\section{Herramientas de medición}\label{sec:c11-medir}

Además del chequeo de reglas hay tres scripts que miden o comprueban cosas puntuales.

\paragraph{\archivo{cruces_capas.py}.} Para cada pista de \capa{In2.Cu} y \capa{B.Cu} lista qué cobre de
\capa{F.Cu} tiene encima, visto desde arriba. Hay una lista de nodos protegidos: los dos nodos \net{SUM}, el nodo
filtrado del DAC, la cadena de \SI{1}{\mega\ohm} y el nodo medio de la realimentación. Son nodos de alta impedancia:
una pista que pase debajo les inyecta ruido por acople capacitivo. Se considera ``encima'' el cobre a menos de
(ancho/2 + \SI{0.2}{\milli\meter}) del eje de la pista. Resultado: \textbf{0 cruces bajo nodos protegidos}. Lo que sí
queda encima de alguna pista interna es IEPE, \net{REST_SV}, SAI e I2C, todas de baja impedancia.

\paragraph{\archivo{conexion_solida.py}.} El datasheet del LMR36015 pide que la tierra de los condensadores de
entrada, de alta frecuencia y de salida y los pines PGND queden unidos por cobre sólido. Un \emph{alivio térmico}
(los ``rayitos'' que unen un pad a una zona) añade inductancia justo en el lazo de corriente más rápido de la placa.
El script comprueba que, en cada buck, los 10 pads de tierra (CIN, CHF, COUT y PGND 1/11) caen enteros en la
\emph{misma} isla del relleno de \capa{F.Cu}, sin alivios. Resultado: cumple en U6 y U7.

\paragraph{\archivo{medir.py}.} Mide caminos de cobre entre pads (solo por \capa{F.Cu}, para que el plano no cuente
como atajo), largos de redes y lazos. Se usó para el lazo de potencia de los bucks, que es lo más sensible al
layout de una fuente conmutada:

\begin{table}[htbp]
\centering
\caption{Mediciones de los bucks con \archivo{medir.py} (iguales en U6 y U7) \medido.}\label{tab:c11-medir}
\small
\begin{tabular}{@{}lr@{}}
\toprule
Medición & Valor\\\midrule
PGND pin 1 $\to$ C29.2 (CHF \SI{220}{\nano\farad}) por \capa{F.Cu} & \SI{2.61}{\milli\meter}\\
PGND pin 11 $\to$ C30.2 por \capa{F.Cu} & \SI{2.21}{\milli\meter}\\
Lazo VIN--CHF--PGND, lado pin 2 & \SI{7.25}{\milli\meter} de perímetro, \SI{3.32}{\milli\meter\squared}\\
Lazo VIN--CHF--PGND, lado pin 10 & \SI{6.97}{\milli\meter}, \SI{2.69}{\milli\meter\squared}\\
Nodo FB (cobre total) & \SI{5.49}{\milli\meter}\\
\bottomrule
\end{tabular}
\end{table}

\section{La verificación después de cada bloque}\label{sec:c11-bloque}

El ruteo se hizo por bloques (capítulo~\ref{cap:ruteo}). Al final de cada uno se corrió
el script \path{tools/ruteo/verificar_bloque.py}, que hace siempre lo mismo y en el mismo orden:
\begin{enumerate}
  \item DRC con paridad, todas las severidades, rellenando y guardando las zonas;
  \item ERC del esquemático;
  \item \archivo{revisar_pistas.py} sobre la placa real;
  \item imágenes de cada capa, para mirarlas.
\end{enumerate}
Solo si todo daba lo esperado se hacía commit y se pasaba al bloque siguiente. Las conexiones sin rutear bajaron de
bloque en bloque (209 al terminar el bloque 1, 0 al terminar el bloque 9). Lo que no daba igual se corregía, o se
presentaba al dueño como decisión antes de seguir.

\begin{idea}[por qué verificar en cada bloque y no al final]
Si un error del bloque 2 se descubre en el bloque 8, hay que deshacer seis bloques de trabajo apoyados encima. Revisar
en cada parada mantiene los errores pequeños y baratos.
\end{idea}

\section{Revisión de los Gerber}\label{sec:c11-gerber}

KiCad no es lo que fabrica la placa: JLCPCB fabrica a partir de los \textbf{Gerber} (un archivo por capa) y de los
archivos de taladro \textbf{Excellon}. Un error al exportar (una capa que falta, unidades mal puestas, un taladro
corrido) no aparece en el DRC de KiCad. Por eso se revisaron los archivos exportados con
\textbf{gerbonara}, una biblioteca de Python que ya existe y lee Gerber. No se escribió un lector propio: un lector
hecho en casa podría tener el mismo error que el exportador y ``confirmarlo''. Se instaló en un entorno aparte con
permiso del dueño. El script es \archivo{tools/revisar_gerber.py} y su informe \archivo{fabricacion/revision_gerber.md}.

\begin{table}[htbp]
\centering
\caption{Resultado de la revisión de los Gerber con gerbonara.}\label{tab:c11-gerber}
\small
\begin{tabular}{@{}L{0.34\linewidth}L{0.58\linewidth}@{}}
\toprule
Comprobación & Resultado\\\midrule
Capas reconocidas & 4 de cobre (F, In1, In2, B), máscaras, serigrafías, pastas, borde y taladro PTH\\
Borde & 4 segmentos, caja de \SI{90.00}{} $\times$ \SI{64.00}{\milli\meter}, contorno cerrado\\
Agujeros & 329 de \SI{0.3}{\milli\meter} (vías), 45 de \SI{1.0}{\milli\meter}, 6 de \SI{1.3}{\milli\meter} y 4 de \SI{2.2}{\milli\meter}; todos PTH\\
Cada PTH sobre un pad en \capa{F.Cu} y \capa{B.Cu} & todos; anillo mínimo \SI{0.150}{\milli\meter}\\
Máscara sobre las vías & 0 aperturas sobre las 329 vías, en las dos caras\\
Serigrafía sobre aperturas de máscara & 0 trazos en las dos caras\\
NPTH & el archivo está vacío (la placa no tiene agujeros sin metalizar); se envía igual\\
\bottomrule
\end{tabular}
\end{table}

\begin{ojo}[un falso positivo que hubo que entender]
La primera vez, la revisión de serigrafía contó cientos de objetos dentro de las aperturas de máscara. No eran
serigrafía: con la opción \texttt{--subtract-soldermask}, KiCad agrega al Gerber de serigrafía unas figuras de
``polaridad clara'' que \emph{borran} la tinta donde hay pads. Se corrigió el script para saltar esos objetos
(348 arriba y 55 abajo) y quedaron 0 trazos reales.
\end{ojo}

\paragraph{Vías tapadas (tenting).} Las vías deben quedar cubiertas, sin apertura en la máscara. Para comprobar que
el método de verificación funciona, se exportó una copia de control con todas las vías destapadas:
\begin{center}
\small
\begin{tabular}{@{}lrrrr@{}}
\toprule
Cara & Aperturas (placa real) & en vías & Aperturas (copia sin tapar) & en vías\\\midrule
\capa{F.Mask} & 348 & \textbf{0} & 677 & 329\\
\capa{B.Mask} & 55 & \textbf{0} & 384 & 329\\
\bottomrule
\end{tabular}
\end{center}
La copia de control encuentra las 329 vías; la placa real no tiene ninguna destapada. Sin la copia de control, un
``0'' podría significar simplemente que el método no ve nada.

\section{El archivo de posiciones (CPL) y el pin 1}\label{sec:c11-cpl}

Para el montaje JLCPCB necesita saber dónde va cada pieza y con qué giro. Eso es el \textbf{CPL}
(\emph{Component Placement List}): designador, posición X e Y, cara y rotación. El problema es que el footprint de
KiCad y el footprint que usa JLC para la misma pieza pueden tener el origen en otro punto o estar dibujados girados.
Si se manda la rotación de KiCad tal cual, la máquina puede poner un integrado girado \SI{90}{\degree}.

\begin{calculo}[corrección de rotación y posición]
\archivo{tools/auditoria_footprints.py} baja el footprint de EasyEDA (el que usa JLC) de cada código LCSC, empareja
los pads por número y busca el giro $d \in \{0, 90, 180, 270\}$ y la traslación $t$ que mejor llevan los pads de KiCad
($p$) sobre los de JLC ($q$):
\begin{equation}
q \approx R(d)\,p + t.
\end{equation}
Con eso, \archivo{tools/fabricacion.py} escribe el CPL con
\begin{equation}
\theta_\text{JLC} = \theta_\text{KiCad} - d,\qquad P_\text{JLC} = P_\text{KiCad} - R(\theta_\text{JLC})\,t.
\end{equation}
Ejemplo: U6 (LMR36015, encapsulado RNX) está a \SI{0}{\degree} en KiCad, pero su footprint de JLC está girado
$d = \SI{90}{\degree}$ y corrido $t = (-0{,}14;\ 0)\,\si{\milli\meter}$. En el CPL va a
\SI{270}{\degree} (es decir, $0 - 90$) con la posición corregida.
\end{calculo}

Para comprobar la cuenta, \archivo{fabricacion/rotaciones.md} coloca el footprint de JLC con la posición y rotación
del CPL y mide la distancia entre su pad 1 y el pad 1 de la placa. En todas las piezas da \textbf{como mucho
\SI{0.21}{\milli\meter}} (los diodos D1 y D2; en los integrados, \SIrange{0}{0.10}{\milli\meter}).

\begin{ojo}[la última palabra la tiene la vista previa de JLC]
La cuenta compara contra el footprint de EasyEDA de cada pieza, pero la colocación real la decide la biblioteca de
montaje de JLC. Por eso, antes de pagar, hay que mirar en la vista previa del pedido: el pin 1 de U1--U10, el cátodo
de D1/D2 del lado de la franja, el pin 1 de Y1 y el centrado de U6/U7 (capítulo~\ref{cap:banco}).
\end{ojo}

\section{Stock de las piezas}\label{sec:c11-stock}

\archivo{tools/stock_bom.py} consulta la API pública de piezas de JLCPCB (solo envía el código LCSC) y para cada
línea del BOM comprueba que el código coincide con el aprobado en \archivo{target.json}, que el número de parte del
fabricante coincide y que hay stock para 5 placas. El 2026-11-16 dio \textbf{35 de 36 líneas bien}. La excepción fue
L2 (XAL4040-103MEC, C5144883) con solo 4 unidades: por eso el dueño decidió ensamblar como máximo 4 placas. El stock
cambia a diario, así que hay que repetir la consulta antes de pedir.

\section{La carpeta de fabricación}\label{sec:c11-fab}

\archivo{tools/fabricacion.py} genera todo lo que se manda a JLCPCB, desde la placa ya verificada, y anota el commit
del que salió (los archivos actuales salieron del commit 970905d; la placa no cambió después). Contenido de
\archivo{hardware/pcb/fabricacion/}:

\begin{longtable}{@{}L{0.33\linewidth}L{0.60\linewidth}@{}}
\caption{Contenido de la carpeta \archivo{fabricacion/}.}\label{tab:c11-fab}\\
\toprule
Archivo & Para qué sirve\\\midrule\endfirsthead
\toprule
Archivo & Para qué sirve\\\midrule\endhead
\archivo{sonometro_revB_gerber.zip} & \textbf{Lo que se sube para fabricar la PCB}: 14 archivos (4 capas de cobre, máscaras, serigrafías, pastas, borde, taladros PTH y NPTH y el archivo de trabajo)\\
\archivo{gerber/} & los mismos archivos sin comprimir, más los mapas de taladro en PDF\\
\archivo{bom.csv} & \textbf{Lista de materiales para el montaje}: Comment, Designator, Footprint, LCSC Part Number. 36 líneas, 112 piezas (solo lo que monta JLC)\\
\archivo{cpl.csv} & \textbf{Posiciones para el montaje}: Designator, Mid X, Mid Y, Layer, Rotation, con la corrección de la sección~\ref{sec:c11-cpl}\\
\archivo{rotaciones.md} & la corrección pieza por pieza y la comprobación del pin 1\\
\archivo{soldar_a_mano.md} & las 8 piezas que no monta JLC (zócalos de la NUCLEO, borneras, BNC, tiras de pines)\\
\archivo{montaje.pdf} & la capa F.Fab con todos los designadores, incluidos los 21 que no caben en la serigrafía\\
\archivo{stock_bom.md} & la consulta de stock del día\\
\archivo{revision_gerber.md} y las vistas SVG & la revisión con gerbonara\\
\archivo{auditoria_piezas.md} & la auditoría de piezas y footprints (capítulo~\ref{cap:piezas})\\
\archivo{LEEME.md} & parámetros del pedido (4 capas, apilado, 1,6 mm, verde, HASL sin plomo, vías taponadas) y la lista de comprobaciones antes de pagar\\
\bottomrule
\end{longtable}

\begin{idea}[qué se sube y dónde]
En la página de pedido de JLCPCB:
\begin{enumerate}
  \item En \emph{PCB} se sube \archivo{sonometro_revB_gerber.zip} y se eligen las opciones de \archivo{LEEME.md}.
  \item Se activa \emph{PCB Assembly} (Economic, cara superior) y, cuando lo pide, se suben \archivo{bom.csv} y
  \archivo{cpl.csv}.
  \item Se revisa la vista previa (pin 1, diodos, Y1, bucks) antes de pagar.
\end{enumerate}
\end{idea}

\section{Estado verificado}\label{sec:c11-estado}

\begin{longtable}{@{}L{0.42\linewidth}L{0.24\linewidth}L{0.26\linewidth}@{}}
\caption{Estado final de las verificaciones (placa del commit 970905d en adelante).}\label{tab:c11-estado}\\
\toprule
Verificación & Resultado & Herramienta\\\midrule\endfirsthead
\toprule
Verificación & Resultado & Herramienta\\\midrule\endhead
DRC, todas las reglas activas, ninguna excluida & 0 violaciones & \texttt{kicad-cli pcb drc}\\
Paridad placa--esquemático & 0 diferencias & \texttt{--schematic-parity}\\
ERC & 0 & \texttt{kicad-cli sch erc}\\
Esquemático contra el diseño & 127 piezas, 68 redes, 0 problemas & \archivo{verify.py}\\
Reglas propias A--M & 0 hallazgos & \archivo{revisar_pistas.py}\\
Conexiones sin rutear & 0 & DRC\\
Cruces bajo nodos protegidos & 0 & \archivo{cruces_capas.py}\\
Tierras de los bucks sin alivios & cumple (U6, U7) & \archivo{conexion_solida.py}\\
Gerber: capas, borde, taladros, anillo & correcto; anillo mín. \SI{0.15}{\milli\meter} & \archivo{revisar_gerber.py}\\
Vías tapadas & 0 de 329 destapadas & Gerber de máscara + copia de control\\
CPL: pin 1 & $\le\SI{0.21}{\milli\meter}$ en todas & \archivo{fabricacion.py}\\
Stock & 35 de 36 líneas (L2: 4 unidades) & \archivo{stock_bom.py}\\
\bottomrule
\end{longtable}

\section*{Resumen}
\begin{itemize}
  \item ERC, paridad y DRC responden preguntas distintas; hacen falta las tres, más un chequeo propio para las reglas
  de buen diseño que KiCad no conoce.
  \item Ninguna regla se ignoró ni se excluyó; los avisos que aparecieron al activarlas se arreglaron de fondo.
  \item \archivo{revisar_pistas.py} tiene 13 chequeos (A--M; el H era de la época de 2 capas) y terminó con 0 hallazgos; las excepciones aceptadas
  quedan listadas aparte con su motivo.
  \item Los Gerber se revisaron con una biblioteca existente y una copia de control para el tenting.
  \item El CPL corrige el giro y el corrimiento de cada footprint respecto del de JLC; el pin 1 queda a
  $\le\SI{0.21}{\milli\meter}$, y se confirma en la vista previa del pedido.
  \item A JLCPCB se sube el zip de Gerber para la PCB y \archivo{bom.csv} + \archivo{cpl.csv} para el montaje.
\end{itemize}
